% !TEX root = ../preview.tex

\section{Preliminaries and Notations}
\label{sec:preliminaries}

% This file mirrors 2_Preliminaries_and_Notations.md. Keep both versions aligned.

\subsection{Problem Formulation: Agent--Environment Interaction}
\label{sec:preliminaries:problem}

We consider an embodied agent interacting with an environment over a family of tasks
$\mathcal{T}$. The interaction is modeled as a controlled Markov process
\begin{equation}
    \mathcal{M} = (\mathcal{S}, \mathcal{A}, P),
    \label{eq:controlled-markov-process}
\end{equation}
where $\mathcal{S}$ and $\mathcal{A}$ denote the environment state space and the
action space, respectively, and $P(s_{t+1}\mid s_t,a_t)$ is the environment
transition kernel. Each task $T\in\mathcal{T}$ may specify an initial-state
distribution $\rho_T$ and a task condition $c_T$, such as a language instruction
or a goal specification. We assume that the tasks share the same state space, action
space, and physical dynamics. At environment time step $t$, the environment is in
state $s_t\in\mathcal{S}$, the agent applies an action $a_t\in\mathcal{A}$, and
the next state evolves according to
\begin{equation}
    s_0\sim\rho_T,
    \qquad
    s_{t+1}\sim P(\cdot\mid s_t,a_t).
    \label{eq:environment-transition}
\end{equation}

The physical state $s_t$ is not necessarily exposed directly to the agent. Let
$\mathcal{O}$ denote the sensing process and $E$ a perceptual encoder. We define
the latent state available to the world--action model (WAM) as
\begin{equation}
    z_t = E\!\left(\mathcal{O}(s_t)\right),
    \label{eq:latent-state-encoding}
\end{equation}
where $z_t\in\mathcal{Z}$ summarizes the sensory observation of $s_t$. This
distinction allows the environment dynamics to remain Markovian in $s_t$, while
the agent may act from encoded and potentially partial observations. We denote the
interaction history available at time $t$ by
\begin{equation}
    \mathcal{H}_t = (c_T,z_0,a_0,z_1,a_1,\ldots,a_{t-1},z_t),
    \label{eq:interaction-history}
\end{equation}
and write the action-selection process generally as
$a_t\sim\pi(\cdot\mid\mathcal{H}_t)$. This notation does not require $z_t$ itself
to be a sufficient Markov state.

Rather than selecting only a single action, a WAM predicts future latent states and
actions over a horizon $H$. Starting at time $t$, we denote its predicted latent
and action chunks by
\begin{equation}
    \hat{\mathbf{Z}}_{t,H}
    = (\hat z_{t+1},\ldots,\hat z_{t+H}),
    \qquad
    \hat{\mathbf{A}}_{t,H}
    = (\hat a_t,\ldots,\hat a_{t+H-1}),
    \label{eq:wam-prediction-chunks}
\end{equation}
with the generic predictive model
\begin{equation}
    (\hat{\mathbf{Z}}_{t,H},\hat{\mathbf{A}}_{t,H})
    \sim p_\theta(\cdot\mid\mathcal{H}_t).
    \label{eq:wam-predictive-model}
\end{equation}

This expression does not prescribe how the two chunks are generated: a WAM may
model them with a joint generative process or factorize generation into successive
latent-state and action stages. During execution, realized transitions and external
disturbances can cause the newly encoded latent $z_{t+1}$ to deviate from its
prediction $\hat z_{t+1}$. Our objective is therefore not to learn a
reward-maximizing policy, but to incorporate such newly available state feedback,
together with previously generated action constraints, into the inference process of
a pretrained WAM. Accordingly, rewards, returns, and value functions are not part of
our formulation.

\subsection{Flow Matching}
\label{sec:preliminaries:flow-matching}

We distinguish an individual variable at one environment step from a chunk spanning
multiple steps. An action
$a_t\in\mathbb{R}^{d_a}$ is a single control command at environment time $t$,
whereas
\begin{equation}
    \mathbf{A}_t
    = [a_t^{\mathsf T},a_{t+1}^{\mathsf T},\ldots,
       a_{t+H-1}^{\mathsf T}]^{\mathsf T}
    \in\mathbb{R}^{D_A},
    \qquad D_A=H d_a,
    \label{eq:action-chunk}
\end{equation}
is the time-stacked action vector of horizon $H$. Similarly,
$z_t\in\mathcal{Z}\subseteq\mathbb R^{d_z}$ denotes the single-step latent state
defined above, while
\begin{equation}
    \mathbf{Z}_t
    = [z_{t+1}^{\mathsf T},z_{t+2}^{\mathsf T},\ldots,
       z_{t+H}^{\mathsf T}]^{\mathsf T}
    \in\mathbb R^{D_Z},
    \qquad D_Z=H d_z,
    \label{eq:latent-chunk}
\end{equation}
denotes the time-stacked future latent-state vector aligned with the action horizon.
Bold uppercase symbols therefore denote chunk-level vectors throughout the paper.
We use
lowercase $x$ for a generic single-step variable and uppercase
$\mathbf{X}$ for its chunk-level counterpart. Thus,
$(x,\mathbf{X})=(a,\mathbf{A})$ for action generation and
$(x,\mathbf{X})=(z,\mathbf{Z})$ for latent-state generation. For a WAM that
generates both modalities jointly, $\mathbf{X}_t$ may instead denote their
concatenation,
\begin{equation}
    \mathbf{X}_t
    =
    \begin{bmatrix}
        \mathbf Z_t\\
        \mathbf A_t
    \end{bmatrix}
    \in\mathbb R^{D_X},
    \qquad D_X=D_Z+D_A.
    \label{eq:joint-wam-chunk}
\end{equation}
For bold vectors $\mathbf x\in\mathbb R^{d_x}$ and
$\mathbf y\in\mathbb R^{d_y}$, all derivatives use the numerator-layout
Jacobian convention,
$\partial\mathbf y/\partial\mathbf x\in\mathbb R^{d_y\times d_x}$.
When $\mathbf X$ denotes a generic generation vector, $D$ denotes its
corresponding dimension.

This generic notation does not imply a particular WAM factorization. A joint WAM
transports the concatenated $\mathbf{X}_t$ with one flow, whereas a stage-wise WAM
applies separate flows to $\mathbf{Z}_t$ and $\mathbf{A}_t$. In either case, the
elements of a generated chunk remain aligned with individual environment steps; this
alignment will later allow FBFM to constrain selected $z_{t+i}$ and
$a_{t+i}$, rather than treating the chunk as an indivisible unit.

Flow Matching learns a time-dependent vector field that transports samples from a
simple source distribution to the conditional data distribution. Let
$\tau\in[0,1]$ denote continuous flow time, distinct from the environment index
$t$. Given a target chunk $\mathbf{X}_t$ and Gaussian noise
$\boldsymbol{\epsilon}\sim\mathcal{N}(\mathbf{0},\mathbf{I})$ of the same dimension,
we use the linear conditional probability path
\begin{equation}
    \mathbf{X}_t^\tau
    = (1-\tau)\boldsymbol{\epsilon}+\tau\mathbf{X}_t,
    \qquad
    \mathbf{X}_t^0=\boldsymbol{\epsilon},
    \quad
    \mathbf{X}_t^1=\mathbf{X}_t.
    \label{eq:linear-flow-path}
\end{equation}
Its conditional target velocity is constant along the path:
\begin{equation}
    \mathbf{u}(\mathbf{X}_t^\tau\mid\mathbf{X}_t)
    = \frac{\mathrm{d}\mathbf{X}_t^\tau}{\mathrm{d}\tau}
    = \mathbf{X}_t-\boldsymbol{\epsilon}.
    \label{eq:conditional-flow-velocity}
\end{equation}

Conditioned on the interaction history $\mathcal{H}_t$, a Flow-Matching model
$v_\theta$ takes the current noisy chunk and flow time as inputs and predicts a
velocity with the same dimension as $\mathbf{X}_t$. It is trained using
\begin{equation}
    \mathcal{L}_{\mathrm{FM}}(\theta)
    =
    \mathbb{E}_{\substack{
        (\mathcal{H}_t,\mathbf{X}_t)\sim\mathcal{D},\,
        \boldsymbol{\epsilon}\sim\mathcal{N}(\mathbf{0},\mathbf{I}),\,
        \tau\sim p(\tau)}}
    \left[
        \left\|
        v_\theta(\mathbf{X}_t^\tau,\tau;\mathcal{H}_t)
        -
        (\mathbf{X}_t-\boldsymbol{\epsilon})
        \right\|_2^2
    \right],
    \label{eq:flow-matching-objective}
\end{equation}
where $p(\tau)$ is a chosen flow-time sampling distribution. Uniform sampling is
the standard choice, while non-uniform schedules may emphasize particular noise
levels.

At inference time, generation starts from
$\mathbf{X}_t^0\sim\mathcal{N}(\mathbf{0},\mathbf{I})$ and solves the conditional
ordinary differential equation
\begin{equation}
    \frac{\mathrm{d}\mathbf{X}_t^\tau}{\mathrm{d}\tau}
    = v_\theta(\mathbf{X}_t^\tau,\tau;\mathcal{H}_t),
    \qquad
    \hat{\mathbf{X}}_t
    = \mathbf{X}_t^0
    + \int_0^1
    v_\theta(\mathbf{X}_t^\tau,\tau;\mathcal{H}_t)\,
    \mathrm{d}\tau.
    \label{eq:flow-matching-ode}
\end{equation}
For example, a forward Euler solver uses
\begin{equation}
    \mathbf{X}_t^{\tau_{k+1}}
    =
    \mathbf{X}_t^{\tau_k}
    + \Delta\tau_k\,
    v_\theta(\mathbf{X}_t^{\tau_k},\tau_k;\mathcal{H}_t).
    \label{eq:flow-matching-euler}
\end{equation}

Some implementations parameterize the same path by the remaining noise level
$\sigma=1-\tau$, which is integrated from $1$ to $0$. Defining the
corresponding field as
$\tilde v_\theta=-v_\theta$, the clean endpoint estimate under the linear path is
\begin{equation}
    \hat{\mathbf{X}}_t^1
    =
    \mathbf{X}_t^\sigma
    -\sigma\,
    \tilde v_\theta(\mathbf{X}_t^\sigma,\sigma;\mathcal{H}_t).
    \label{eq:clean-endpoint-estimate}
\end{equation}
We use superscripts exclusively for flow time or noise level and subscripts for
environment time. This separation is important for FBFM, which modifies the
inference-time vector field while preserving the pretrained Flow-Matching model.

For quick reference, Appendix~\ref{app:notation-index} provides a consolidated
index of the notation used throughout the paper.
